% GENERATED FILE. Edit canonical lessons, then run npm run generate:books.
% canonical-source-hash: fnv1a64:6273213a5bcec629
% canonical-lessons: ML-C431-vikasanam, ML-C431-vikasippikkuka, ML-C431-navikarikkuka, ML-C431-purogamikkuka, ML-C431-sampradayam

\chapter{Development}
\label{ch:ml-development}
\hlchaptermodality{431}

\begin{chapteropening}
\textbf{By the end of this chapter:} \emph{I can say development, to develop, to modernise, to progress and a system.}

Complete the last lesson of chapter 431: I can say development, to develop, to modernise, to progress and a system.
\end{chapteropening}

\section[\texorpdfstring{\ml{വികസനം} (vikasanaṁ)}{വികസനം (vikasanaṁ)}]{\ml{വികസനം} (vikasanaṁ) — development}

\label{lesson:ML-C431-vikasanam}



\begin{quote}
Before the new one: say the Malayalam for transparent, then the Malayalam for to investigate.
\end{quote}

\subsection*{You'll want to know: \ml{വികസനം}}
\textbf{\ml{വികസനം}} — \emph{vikasanaṁ} — \textquotedblleft{}development\textquotedblright{}.

\begin{quote}
\ml{ഗ്രാമത്തിന്റെ} \ml{വികസനത്തിന്} \ml{പണം} \ml{വേണം}.
\end{quote}

\begin{quote}
\emph{grāmattinṟe vikasanattinŭ paṇaṁ vēṇaṁ.}
\end{quote}

\begin{quote}
The village needs money for its development.
\end{quote}

Of a village, a town or a country: \ml{ഗ്രാമവികസനം}, rural development.

\begin{cousinweb}
\textbf{\ml{വികസനം}} (\emph{vikasanaṁ}) is built on the Sanskrit root \emph{vi-kas}, \textquotedblleft{}to open, to expand\textquotedblright{}, whose noun \emph{vikāsa} means \textquotedblleft{}budding; expansion, growth\textquotedblright{} (Monier-Williams). Development is a flower opening.
\end{cousinweb}

\begin{grammarlens}[title={What you've built}]
One more word for giving reasons and opinions.
\end{grammarlens}

\subsection*{Your turn}
\begin{itemize}
\raggedright
  \item \emph{Say it:} \emph{vikasanaṁ}
  \item \emph{Say it:} \emph{vikasanaṁ}, once more
  \item \emph{Say it:} \emph{grāmattinṟe vikasanattinŭ paṇaṁ vēṇaṁ.} — The village needs money for its development.
  \item \emph{Recall:} say the Malayalam for transparent, then the Malayalam for to investigate, then say \emph{vikasanaṁ} again
\end{itemize}

\begin{tcolorbox}[breakable,colback=teal!4,colframe=teal!35!black,title={Before you move on}]
What does \ml{വികസനം} mean? (\textquotedblleft{}Development\textquotedblright{}.) Say it once more.
\end{tcolorbox}
\section[\texorpdfstring{\ml{വികസിപ്പിക്കുക} (vikasippikkuka)}{വികസിപ്പിക്കുക (vikasippikkuka)}]{\ml{വികസിപ്പിക്കുക} (vikasippikkuka) — to develop}

\label{lesson:ML-C431-vikasippikkuka}



\begin{quote}
Before the new one: say the Malayalam for to investigate, then the Malayalam for development.
\end{quote}

\subsection*{You'll want to know: \ml{വികസിപ്പിക്കുക}}
\textbf{\ml{വികസിപ്പിക്കുക}} — \emph{vikasippikkuka} — \textquotedblleft{}to develop\textquotedblright{}.

\begin{quote}
\ml{അവർ} \ml{ഒരു} \ml{പുതിയ} \ml{ആപ്പ്} \ml{വികസിപ്പിച്ചു}.
\end{quote}

\begin{quote}
\emph{avar oru putiya āppŭ vikasippiccu.}
\end{quote}

\begin{quote}
They developed a new app.
\end{quote}

Past: \ml{വികസിപ്പിച്ചു}, \emph{vikasippiccu}.

Its intransitive partner is \ml{വികസിക്കുക}, to grow, to expand.

\begin{cousinweb}
\textbf{\ml{വികസിപ്പിക്കുക}} (\emph{vikasippikkuka}) is built from the Sanskrit root \emph{vi-kas}, \textquotedblleft{}to open, to expand\textquotedblright{} (Monier-Williams), with the Malayalam causative verb ending -\ml{ിപ്പിക്കുക}: to make a thing grow.
\end{cousinweb}

\begin{grammarlens}[title={What you've built}]
One more word for giving reasons and opinions.
\end{grammarlens}

\subsection*{Your turn}
\begin{itemize}
\raggedright
  \item \emph{Say it:} \emph{vikasippikkuka}
  \item \emph{Say it:} \emph{vikasippikkuka}, once more
  \item \emph{Say it:} \emph{avar oru putiya āppŭ vikasippiccu.} — They developed a new app.
  \item \emph{Recall:} say the Malayalam for to investigate, then the Malayalam for development, then say \emph{vikasippikkuka} again
\end{itemize}

\begin{tcolorbox}[breakable,colback=teal!4,colframe=teal!35!black,title={Before you move on}]
What does \ml{വികസിപ്പിക്കുക} mean? (\textquotedblleft{}To develop\textquotedblright{}.) Say it once more.
\end{tcolorbox}
\section[\texorpdfstring{\ml{നവീകരിക്കുക} (navīkarikkuka)}{നവീകരിക്കുക (navīkarikkuka)}]{\ml{നവീകരിക്കുക} (navīkarikkuka) — to modernise; to renovate}

\label{lesson:ML-C431-navikarikkuka}



\begin{quote}
Before the new one: say the Malayalam for development, then the Malayalam for to develop.
\end{quote}

\subsection*{You'll want to know: \ml{നവീകരിക്കുക}}
\textbf{\ml{നവീകരിക്കുക}} — \emph{navīkarikkuka} — \textquotedblleft{}to modernise; to renovate\textquotedblright{}.

\begin{quote}
\ml{പഞ്ചായത്ത്} \ml{പഴയ} \ml{ലൈബ്രറി} \ml{നവീകരിച്ചു}.
\end{quote}

\begin{quote}
\emph{pañcāyattŭ paḻaya laibraṟi navīkariccu.}
\end{quote}

\begin{quote}
The panchayat renovated the old library.
\end{quote}

Past: \ml{നവീകരിച്ചു}, \emph{navīkariccu}.

The noun is \ml{നവീകരണം}, renovation.

\begin{cousinweb}
\textbf{\ml{നവീകരിക്കുക}} (\emph{navīkarikkuka}) is built from Sanskrit \emph{nava}, \textquotedblleft{}new, young\textquotedblright{} (CDIAL 6983), with -\ml{ീകരിക്കുക}, \textquotedblleft{}to make into\textquotedblright{}: to make a thing new.
\end{cousinweb}

\begin{grammarlens}[title={What you've built}]
One more word for giving reasons and opinions.
\end{grammarlens}

\subsection*{Your turn}
\begin{itemize}
\raggedright
  \item \emph{Say it:} \emph{navīkarikkuka}
  \item \emph{Say it:} \emph{navīkarikkuka}, once more
  \item \emph{Say it:} \emph{pañcāyattŭ paḻaya laibraṟi navīkariccu.} — The panchayat renovated the old library.
  \item \emph{Recall:} say the Malayalam for development, then the Malayalam for to develop, then say \emph{navīkarikkuka} again
\end{itemize}

\begin{tcolorbox}[breakable,colback=teal!4,colframe=teal!35!black,title={Before you move on}]
What does \ml{നവീകരിക്കുക} mean? (\textquotedblleft{}To modernise; to renovate\textquotedblright{}.) Say it once more.
\end{tcolorbox}
\section[\texorpdfstring{\ml{പുരോഗമിക്കുക} (purōgamikkuka)}{പുരോഗമിക്കുക (purōgamikkuka)}]{\ml{പുരോഗമിക്കുക} (purōgamikkuka) — to make headway; to progress}

\label{lesson:ML-C431-purogamikkuka}



\begin{quote}
Before the new one: say the Malayalam for to develop, then the Malayalam for to modernise.
\end{quote}

\subsection*{You'll want to know: \ml{പുരോഗമിക്കുക}}
\textbf{\ml{പുരോഗമിക്കുക}} — \emph{purōgamikkuka} — \textquotedblleft{}to make headway; to progress\textquotedblright{}.

\begin{quote}
\ml{പണി} \ml{വേഗത്തിൽ} \ml{പുരോഗമിക്കുന്നു}.
\end{quote}

\begin{quote}
\emph{paṇi vēgattil purōgamikkunnu.}
\end{quote}

\begin{quote}
The work is progressing fast.
\end{quote}

Past: \ml{പുരോഗമിച്ചു}, \emph{purōgamiccu}.

Said of work, talks or a plan, not of a person walking.

\begin{cousinweb}
\textbf{\ml{പുരോഗമിക്കുക}} (\emph{purōgamikkuka}) is built from Sanskrit \emph{purogama}, \textquotedblleft{}going in front, leading\textquotedblright{} (Monier-Williams), with the Malayalam verb ending -\ml{ിക്കുക}.
\end{cousinweb}

\begin{grammarlens}[title={What you've built}]
One more word for giving reasons and opinions.
\end{grammarlens}

\subsection*{Your turn}
\begin{itemize}
\raggedright
  \item \emph{Say it:} \emph{purōgamikkuka}
  \item \emph{Say it:} \emph{purōgamikkuka}, once more
  \item \emph{Say it:} \emph{paṇi vēgattil purōgamikkunnu.} — The work is progressing fast.
  \item \emph{Recall:} say the Malayalam for to develop, then the Malayalam for to modernise, then say \emph{purōgamikkuka} again
\end{itemize}

\begin{tcolorbox}[breakable,colback=teal!4,colframe=teal!35!black,title={Before you move on}]
What does \ml{പുരോഗമിക്കുക} mean? (\textquotedblleft{}To make headway; to progress\textquotedblright{}.) Say it once more.
\end{tcolorbox}
\section[\texorpdfstring{\ml{സമ്പ്രദായം} (sampradāyaṁ)}{സമ്പ്രദായം (sampradāyaṁ)}]{\ml{സമ്പ്രദായം} (sampradāyaṁ) — a system; a tradition}

\label{lesson:ML-C431-sampradayam}



\begin{quote}
Before the new one: say the Malayalam for to modernise, then the Malayalam for to make headway.
\end{quote}

\subsection*{You'll want to know: \ml{സമ്പ്രദായം}}
\textbf{\ml{സമ്പ്രദായം}} — \emph{sampradāyaṁ} — \textquotedblleft{}a system; a tradition\textquotedblright{}.

\begin{quote}
\ml{പഴയ} \ml{സമ്പ്രദായം} \ml{മാറ്റാൻ} \ml{സമയമായി}.
\end{quote}

\begin{quote}
\emph{paḻaya sampradāyaṁ māṟṟān samayamāyi.}
\end{quote}

\begin{quote}
It is time to change the old system.
\end{quote}

Plural: \ml{സമ്പ്രദായങ്ങൾ}, customs.

\begin{cousinweb}
\textbf{\ml{സമ്പ്രദായം}} (\emph{sampradāyaṁ}) is borrowed from Sanskrit \emph{sampradāya}, \textquotedblleft{}tradition, a doctrine handed down from one teacher to the next\textquotedblright{} (Monier-Williams).
\end{cousinweb}

\begin{grammarlens}[title={What you've built}]
One more word for giving reasons and opinions.
\end{grammarlens}

\subsection*{Your turn}
\begin{itemize}
\raggedright
  \item \emph{Say it:} \emph{sampradāyaṁ}
  \item \emph{Say it:} \emph{sampradāyaṁ}, once more
  \item \emph{Say it:} \emph{paḻaya sampradāyaṁ māṟṟān samayamāyi.} — It is time to change the old system.
  \item \emph{Recall:} say the Malayalam for to modernise, then the Malayalam for to make headway, then say \emph{sampradāyaṁ} again
\end{itemize}

\begin{tcolorbox}[breakable,colback=teal!4,colframe=teal!35!black,title={Before you move on}]
What does \ml{സമ്പ്രദായം} mean? (\textquotedblleft{}A system; a tradition\textquotedblright{}.) Say it once more.
\end{tcolorbox}
